\documentclass{article}
\usepackage{multirow}
\usepackage[table,xcdraw]{xcolor}
\usepackage{booktabs} % 表の横線を調整するためのパッケージ
\usepackage{tabularx} % 表の幅を調整するためのパッケージ

\newcommand{\bettershortstack}[4][c]{%
  \renewcommand{\arraystretch}{#2}
  \begin{tabular}[b]{@{}#1@{}}
  #4
  \end{tabular}%
  \renewcommand{\arraystretch}{#3}
}
%\bettershortstack[左右寄せ]{行間の指定}{下の行間}{セル内の文章}


\begin{document}

\begin{table}[htbp]
\centering
\caption{ログイン画面を生成するレイヤー}
\begin{tabularx}{\linewidth}{llp{20cm}}
    
    \cellcolor[HTML]{C0C0C0}\textbf{レイヤー名} & \multicolumn{2}{l}{login\_page.dart} \\ 
    \cellcolor[HTML]{C0C0C0}\textbf{レイヤー概要} & \multicolumn{2}{l}{ログイン画面を生成するレイヤー} \\
    \cellcolor[HTML]{C0C0C0}\textbf{責任者} & \multicolumn{2}{l}{山本昇永} \\ 
    \midrule
    \multicolumn{1}{l}{\cellcolor[HTML]{C0C0C0}\textbf{構成要素}}       & \multicolumn{2}{l}{\cellcolor[HTML]{C0C0C0}}  \\
& \cellcolor[HTML]{EFEFEF}\textbf{クラス名}   & \multicolumn{1}{l}{LoginPage}                \\
    \cmidrule{2-3}
& \cellcolor[HTML]{EFEFEF}\textbf{クラスの種類} & \multicolumn{1}{l}{extends: StatefulWidget}  \\
& \cellcolor[HTML]{EFEFEF}\textbf{処理概要}   & \multicolumn{1}{l}{ログイン画面をWidgetを動的に変更するクラス} \\
\cmidrule{2-3}
& \multicolumn{2}{l}{\cellcolor[HTML]{EFEFEF}\textbf{入力}}                                \\
& \multicolumn{2}{l}{なし}     \\
& \multicolumn{2}{l}{\cellcolor[HTML]{EFEFEF}\textbf{出力}} \\
& \multicolumn{2}{l}{なし} \\
\cmidrule{2-3}
& \multicolumn{2}{l}{\cellcolor[HTML]{EFEFEF}\textbf{他クラス・関数との関係}} \\
& \multicolumn{2}{l}{・login\_page.dartのLoginPageStateクラスを組み込む} \\
\cmidrule{2-3}
& \cellcolor[HTML]{EFEFEF}\textbf{クラス名}   & \multicolumn{1}{l}{LoginPage}                \\
    \cmidrule{2-3}
& \cellcolor[HTML]{EFEFEF}\textbf{クラスの種類} & \multicolumn{1}{l}{extends: StatefulWidget}  \\
& \cellcolor[HTML]{EFEFEF}\textbf{処理概要}   & \multicolumn{1}{l}{ログイン画面をWidgetを動的に変更するクラス} \\
\cmidrule{2-3}
& \multicolumn{2}{l}{\cellcolor[HTML]{EFEFEF}\textbf{入力}}                                \\
& \multicolumn{2}{l}{なし}     \\
& \multicolumn{2}{l}{\cellcolor[HTML]{EFEFEF}\textbf{出力}} \\
& \multicolumn{2}{l}{なし} \\
\cmidrule{2-3}
& \multicolumn{2}{l}{\cellcolor[HTML]{EFEFEF}\textbf{他クラス・関数との関係}} \\
& \multicolumn{2}{l}{・new\_menber\_company.dartのNewMemberGeneralクラスを} \\
& \multicolumn{2}{l}{呼び出すことで、一般用の新規会員登録画面を表示} \\
\bottomrule

    
\end{tabularx}
\end{table}
\end{document}